\scanpage{286}
\mspage{27}

Je vais \deleted{la} désigner cette catégorie fibrée par $\underline{\mathrm{Ch}}\unsure{\cdot|.}(X)$. (Provisoirement -- si elle s'avère trop grosse, je la \unsure{restreindrai|restraindrai}\ldots). Ainsi on a
\diagram{\begin{tikzcd}[ampersand replacement=\&]\underline{\mathrm{Ch}}(X) \arrow[d,"{\pi_X}"]\\ \underline{\mathrm{Ty}}\end{tikzcd}}.

La fibre de $\underline{\mathrm{Ch}}(X)$ en $\tau\in\mathrm{Ob}\,\underline{\mathrm{Ty}}$ est la catégorie $\underline{\mathrm{Ch}}_\tau(X)$ des chemins de type $\tau$ dans $X$. \deleted{Mais \unsure{qu}} Cela nous dit donc quelles sont les flèches de \deleted{\unsure{$X$}}\inserted{$\underline{\mathrm{Ch}}(X)$} au dessus des flèches identiques de $\underline{\mathrm{Ty}}$. Mais que sont les flèches \inserted{plus} générales ? On se donne une flèche dans $\underline{\mathrm{Ty}}$
\deleted{$I_{\tau'}\longrightarrow I_\tau$}
\[ \varphi\colon\tau'\longrightarrow\tau\qquad\text{i.e.}\quad I_{\tau'}\xrightarrow{\varphi}I_\tau\quad\text{dans }\mathrm{Cat}, \]
quelles sont les flèches de $\underline{\mathrm{Ch}}(X)$ au dessus de $\varphi$ ? Par le sorite général des catégories fibrées, si
\[ \begin{aligned}
c&\in\underline{\mathrm{Ch}}(X)_\tau=\underline{\mathrm{Ch}}_\tau(X)\overset{\mathrm{def}}{=}\underline{\mathrm{Hom}}(I_\tau,X)\\
c'&\in\underline{\mathrm{Ch}}(X)_{\tau'}=\underline{\mathrm{Ch}}_{\tau'}(X)\overset{\mathrm{def}}{=}\underline{\mathrm{Hom}}(I_{\tau'},X),
\end{aligned} \]
on a
\[ \boxed{\mathrm{Hom}_\varphi(\unsure{c'|\tau'},c)\simeq\mathrm{Hom}\left(c',\underbrace{\varphi^*(c)}_{\substack{\Vert\\c\circ\varphi}}\right)} \]
Donc \unsure{!|:} les \deleted{flèches \unsure{dans}} homomorphismes d'un
% archive sheet number: 285
